% PLAN:N1
% NODES:
\section{Introduction}\label{sec:N1}

% N1-p1
Questions about a many-body Hamiltonian $H$ probed through a state $|\psi\rangle$ reduce to questions about the probe-weighted spectral measure $\mu_\psi = \sum_j p_j\,\delta(E - E_j)$ - where its weight sits, at which energy it is read, and at what resolution. A quantum device that evolves $|\psi\rangle$ under $H$ accesses this measure through the real-time overlap $G(t) = \langle\psi|e^{-iHt}|\psi\rangle$, its characteristic function, whose frequencies are the spectrum and whose amplitudes are the spectral weights. The operational problem is to turn a finite set of measured overlap values into spectral statements at chosen locations and resolutions, at stated accuracy - to read regularised spectral quantities at a chosen location and scale, not to recover individual eigenvalues.

% N1-p2
Conventional quantum spectral protocols implement the spectral question on the device. Phase-estimation-type readouts, coherently synthesised filters and subspace or eigenstate-targeting schemes fix the target energy and the effective resolution in the circuit that runs, so a new location or a new resolution means a new quantum experiment, and the measurement record produced for one question does not in general serve the next. This is no defect of any particular scheme, and Sec.~\ref{sec:N2} sets the present method against published ones task by task. What the designs share is a single assumption, uniform enough to pass unstated - that spectral location and resolution are parameters of acquisition.

% N1-p3
Spectral Surface Generation (SSG) removes this assumption: the spectral filter moves off the device and into classical post-processing. A quantum device samples the spectral signal $G(t)$ at fixed Gauss-Hermite grid times $t_k$ in one acquisition campaign, and the estimates $\hat{G}(t_k)$ are banked. A classical layer then reconstructs the Gaussian-filtered spectral surface at a spectral location $s$ and inverse-scale resolution $\tau$ chosen after acquisition. New $s$, new $\tau$ or a coefficient insertion means only new reconstruction coefficients applied to the stored record - no further quantum evolution. We carry the derivation in the spectral coordinate $s$, the actual spectral location at which the Gaussian window is centred. The one bridge to the parent ITQDE convention is $s = T - \lambda$, with $\lambda$ the displacement from a chosen target $T$, and $\lambda$ appears only there. One narrow claim organises the method: the measured overlaps are independent of the target and support a range of resolutions. Sec.~\ref{sec:N2} proves the independence for the measured objects, and Sec.~\ref{sec:N3} derives the reconstruction identity behind the dials. The independence is a property of the acquisition, not a promise that every target is equally readable - what can be read, at what accuracy, is bounded in Sec.~\ref{sec:N4}.

% N1-p4
The principal object is the spectral surface $Z_\tau(s) = \sum_j p_j\, e^{-\tau (E_j - s)^2}$ - a Gaussian window of width $1/\sqrt{\tau}$, centred at $s$, applied to the probe-weighted spectral measure - so small $\tau$ is a broad survey and large $\tau$ a narrow local readout. From the same banked record follow the generating surface $\Phi_\tau(s) = \log Z_\tau(s)$, whose derivatives are the response channels - the filtered mean, the filtered variance and the higher cumulants - and, beyond them, within-window ratios and observable branches. Figure~\ref{fig:nf1} fixes the architecture before any formalism. Panel (a) is the pipeline $\hat{G}(t_k) \to c_k(s,\tau) \to \hat{Z}_\tau(s)$, with the quantum-classical boundary drawn once and the three causes of the finite-record envelope labelled at their sources. Panel (b) is a three-level toy measure under three Gaussian kernels, its surface values computed as exact spectral sums. The figure is a schematic and an exact toy example - it establishes nothing about finite-shot accuracy, about any specific Hamiltonian family or about resource costs, and the envelope causes it labels are derived only in Sec.~\ref{sec:N4}.

% N1-p5
The reuse is bounded, and the bounds are part of the method's definition. Hats mark finite-record objects: $\hat{Z}_{\tau,m}(s)$ is reconstructed from $m$ retained nodes, while the unhatted $Z_\tau(s)$ is reserved for the exact continuous object and the algebraic identities. One banked record supports a finite $(s, \tau)$ domain, the finite-record envelope, bounded by three causes derived in Sec.~\ref{sec:N4}: time reach from the Gaussian aperture, node retention from the Gauss-Hermite gate with its measured retention law $\bar{m} \sim \|H\|\sqrt{\tau}$, and the shot-noise floor at shot budget $K$ per retained estimator. Shot noise propagates through the linear reconstruction with a derived covariance for estimators sharing one record, and the cost of the classical layer is accounted under an explicitly named query model in Sec.~\ref{sec:N3}. An absolute read of exponentially small surface values is expensive in shots, so the method reads no exponentially small spectral weight at polynomial shot cost, and every later small-weight statement is sized against that covariance and floor.

% N1-p6
We validate the method on a controlled subset-sum family, chosen because it is checkable, not because it is hard. The instance $a = (a_1, \ldots, a_n)$ is encoded in the diagonal Hamiltonian $H_0$ and never in the probe. At $V = 0$ the overlap record collapses exactly to computable Fourier coefficients, so every reconstructed quantity has an independently computable classical reference, and at the production weight scale $a_i = O(n^2)$ dynamic programming produces that reference in polynomial time. The benchmark therefore inherits no worst-case NP-hardness from subset-sum - at this weight scale the instances are pseudo-polynomially tractable - and nothing in the choice of family asserts a separation between the method and classical computation. Bounded noncommuting perturbations $V$ supply the diagnostic regime - a multiplicity fingerprint under a stated forward model, degeneracy cartography and an integer decoder with an erasure rule - constructed in Sec.~\ref{sec:N5}. A planted equal-weight family probes whether readability and growing classical simulation cost can coexist over the tested range, and its interpretation belongs to Sec.~\ref{sec:N7}.

% N1-p7
The paper keeps its exact statements and its finite-size observations apart, and that separation holds throughout. The reconstruction, response and cumulant identities, together with the $V = 0$ coefficient collapse, are exact algebra. The numerical results are finite-size observations at proof-of-principle sizes - exact diagonalisation at $n \le 12$, TEBD matrix product state evolution at $n \le 20$, synthetic finite-shot sampling at stated shot budgets - and the derived guides drawn alongside them for comparison are neither fitted nor asserted as laws. No quantum hardware data appears. The method asserts neither asymptotic hardness nor quantum advantage over classical approaches: finite-size compressibility of the tested generic instances is no hardness no-go, and the reconstruction is validated against a classically computed reference under trusted assumptions rather than certified. The hardness question raised by the planted family stays open, and the bond-dimension growth observed there, consistent with exponential scaling across the tested range, is a guide to that question rather than an established law.

% N1-p8
The argument runs as a dependency chain. Sec.~\ref{sec:N2} reconstructs the parent ITQDE machinery that produces the banked overlap record, proves the record's independence of the target and locates the method task by task against record-based spectral approaches. Sec.~\ref{sec:N3} defines Spectral Surface Generation by its input-output contract and derives the reconstruction, response and cumulant identities that make $s$, $\tau$ and coefficient insertions post-acquisition dials. Sec.~\ref{sec:N4} derives the acquisition chain and bounds the supported $(s, \tau)$ domain of one banked record. Sec.~\ref{sec:N5} instantiates the controlled subset-sum embodiment and defines its validation reference. Sec.~\ref{sec:N6} validates the reconstruction, the fingerprint and the decoder against that reference, distinguishing its exact results from its finite-size observations. Sec.~\ref{sec:N7} bounds what the results establish and states the open hardness question. Sec.~\ref{sec:N8} closes with the three statements the body supports. The appendices carry the complete derivations, the data and code manifest, and the disclosure crosswalk, which records how exactly each supported statement is established.
